\section{Conclusion}\label{sec:concl}

\emph{[To be written once \cref{sec:eval} is in place.]}

% =========================================================================
% TODO --- outline, to be filled once the numbers exist:
%
%   * Restate the framing of Section \ref{sec:intro}: the successor function is
%     the easy half, the passed set is the hard half, and the hard half is hard
%     because its operation is a read-mostly, order-sensitive, expensive update
%     rather than an insert-if-absent.
%
%   * Restate what the paper established, with the evaluation numbers attached:
%     coverage monotonicity (Lemma \ref{lem:monotone}) as the property that lets
%     a shared subsumption test be trusted without validation; termination and
%     verdict determinism without assuming an antichain (Theorems \ref{thm:term}
%     and \ref{thm:verdict}); and the three cost mechanisms, each proved not to
%     change what is explored.
%
%   * The transferable design rule: a compressed representation must be chosen
%     with its query pattern in view, not only with its size, and an adaptive
%     policy should be driven by a signal the engine already computes (visits x
%     size for the mirror, zones per discrete state for the abstraction switch)
%     rather than by a model-specific constant.
%
%   * Limitations, stated plainly: the stored set is not reproducible across
%     runs; strict level synchronisation leaves a tail-idle term that a
%     work-stealing design would not have, at the cost of more stored states;
%     the engine is shared-memory only.
%
%   * Future work: liveness properties, which need a different passed-set
%     discipline; distributed execution over several shared-memory nodes; and
%     using the same monotonicity argument to justify a lock-free stored-set
%     update.
% =========================================================================
